\makeatletter
\def\input@path{{./}{../../../preamble/}{book2026-2027/preamble/}}
\makeatother
\input{preamble.tex}
\usepackage{needspace}
\hypersetup{pdftitle={Grade 10 Term 2 Discrete Probability}}
\setlength{\emergencystretch}{2em}
\title{Grade 10, Term 2, Lesson 3, Criterion C9\\Selecting and Justifying Discrete Probability Models}
\author{}
\date{}
\begin{document}
\maketitle
\setcounter{tocdepth}{1}
\setcounter{secnumdepth}{2}
\phantomsection\label{toc}
\tableofcontents
\section*{Evaluated criterion}
\textbf{C9:} Compares Bernoulli, binomial, negative binomial, geometric, and Poisson distributions by selecting the most appropriate model for a given financial-risk situation.
\section{Introduction}\label{sec:t2-c9-introduction}
All data in this activity are hypothetical teaching models. Probabilities have no monetary units. Keep exact expressions where useful and round final probabilities to four decimal places.
\subsection{Ask what is random}
A school enterprise can record whether one visitor buys, sales among ten visitors, contacts until its first sale, contacts until its third sale, or arrivals during an hour. These are different questions even if all describe the same shop. Start by defining the variable; the word sales alone does not choose a distribution.
\textbf{Check.} One customer either buys or does not: which model? \emph{Suggested answer:} Bernoulli, if 1 labels buying and 0 labels not buying.


\subsection{Find what is fixed}
Fixed trials with random success count suggest binomial. Fixed target successes with random total trials suggest geometric for target one or negative binomial for target $r$. Event counts over fixed exposure suggest Poisson when its process assumptions hold. Counts have count units, probabilities are unitless and monetary consequences need a separate definition.
\textbf{Check.} Ten calls regardless of results: is this geometric? \emph{Suggested answer:} No. A success count over ten fixed calls suggests binomial.


\subsection{Recall notation before choosing}
The symbol $\sim$ is read as has the distribution. In a PMF, $P(X=k)$ is the probability of exactly value $k$. The factorial is $k!=k(k-1)\cdots1$, with $0!=1$. Combinations count success-position choices: $\binom nk=n!/[k!(n-k)!]$. A sum $\sum_{j=0}^{b}$ adds the terms for integers 0 through $b$. In Poisson, $e\approx2.71828$ and $\lambda$ (lambda) is the interval mean, not automatically a rate.
\textbf{Check.} Which model counts trials until success three? \emph{Suggested answer:} Negative binomial with target $r=3$, including the final successful trial.

\textbf{Check.} Is a rare default automatically Poisson? \emph{Suggested answer:} No. Define the exposure and check constant rate, independent interval counts and isolated events.

\section{General worked problems}
\Needspace{14\baselineskip}
\subsection{Introductory general problem: same trials, different variables}
Independent binary trials have constant success probability $1/2$. Variable $A$ indicates success on trial one; $B$ counts successes in three fixed trials; $T$ counts trials until first success.

\textbf{Questions/Tasks}
\begin{enumerate}
\item Select each model, state parameters and support, reject an alternative and calculate $P(A=1)$, $P(B=1)$ and $P(T=3)$.
\end{enumerate}
\subsubsection*{Complete worked solution}
\begin{center}
\begin{tabular}{llp{0.48\linewidth}}
\toprule
Symbol & Value & Meaning in this problem \\ \midrule
$p$ & $0.5$ & Success probability per trial \\
$n$ & $3$ trials & Fixed trials for variable B only \\
\bottomrule
\end{tabular}
\end{center}

\par\noindent\begin{minipage}{\linewidth}
\textbf{Step 1: Define and select.} $A\sim\operatorname{Bernoulli}(0.5)$ with support $\{0,1\}$, $B\sim\operatorname{Bin}(3,0.5)$ with support $\{0,1,2,3\}$ successes, and geometric $T$ with $p=0.5$ and support $\{1,2,\ldots\}$ trials.
\end{minipage}\par\medskip
\par\noindent\begin{minipage}{\linewidth}
\textbf{Step 2: Check and reject.} Binary independent constant-probability trials are supplied. A is one label rather than a trial count; B fixes trials rather than waiting for success; T waits rather than observing all three fixed trials. These distinctions reject the corresponding plausible alternatives.
\end{minipage}\par\medskip
\par\noindent\begin{minipage}{\linewidth}
\textbf{Step 3: Calculate.} $P(A=1)=p=0.5000$; $P(B=1)=\binom31p(1-p)^2=3(0.5)(0.5)^2=0.3750$; $P(T=3)=(1-p)^2p=(0.5)^2(0.5)=0.1250$.
\end{minipage}\par\medskip
\par\noindent\begin{minipage}{\linewidth}
\textbf{Step 4: Check and interpret.} The last event is FFS, whereas the middle event includes SFF, FSF and FFS. The threefold difference checks that the stopping rule was used correctly.
\end{minipage}\par\medskip
\subsection{Bridge to the general model}
\textbf{Comparison reference.} Supports below assume $0<p<1$ and $\lambda>0$. Trials are binary; success must always be stated explicitly. All PMFs are zero outside their stated support.
\begin{center}\small
\begin{tabularx}{\linewidth}{@{}p{0.13\linewidth}p{0.43\linewidth}X@{}}
\toprule Model & What is observed, fixed rule, support and parameters & Assumptions and PMF\\\midrule
Bernoulli & One binary label; one observation in a specified period. Support $\{0,1\}$. Parameter $p$. & Two exhaustive, exclusive labels. No repeated-trial independence required. $P(X=x)=p^x(1-p)^{1-x}$.\\[5pt]
Binomial & Number of successes; fixed $n$ trials. Support $\{0,\ldots,n\}$. Parameters $n,p$. & Independent binary trials with constant $p$. $P(X=k)=\binom nkp^k(1-p)^{n-k}$.\\[5pt]
Geometric & Total trials to first success; stop at success one. Support $\{1,2,\ldots\}$. Parameter $p$. & Independent binary trials with constant $p>0$. $P(T=t)=(1-p)^{t-1}p$.\\[5pt]
Negative binomial & Total trials to success $r$; stop at fixed target $r$. Support $\{r,r+1,\ldots\}$. Parameters $r,p$. & Independent binary trials with constant $p>0$. $P(T=t)=\binom{t-1}{r-1}p^r(1-p)^{t-r}$.\\[5pt]
Poisson & Event count in fixed time, area or exposure. Support $\{0,1,\ldots\}$. Parameter $\lambda=\text{rate}\times\text{exposure}$. & Constant rate, independent disjoint-interval counts, isolated events. $P(X=k)=e^{-\lambda}\lambda^k/k!$.\\
\bottomrule
\end{tabularx}\end{center}
\textbf{Model-selection flowchart.} Follow a yes arrow to a candidate, then check its assumptions. A candidate is not a confirmed model until those checks pass.
\begin{center}
\begin{tikzpicture}[q/.style={draw,rounded corners,text width=7cm,align=center,font=\small},a/.style={draw,text width=5cm,align=center,font=\small},>=Stealth]
\node[q] (q1) at (0,0) {One binary indicator?};
\node[a] (a1) at (8,0) {Bernoulli: identify $p$};
\node[q] (q2) at (0,-1.2) {Success count in fixed $n$ trials?};
\node[a] (a2) at (8,-1.2) {Binomial: identify $n,p$};
\node[q] (q3) at (0,-2.4) {Total trials until first success?};
\node[a] (a3) at (8,-2.4) {Geometric: identify $p$};
\node[q] (q4) at (0,-3.6) {Total trials until success $r$?};
\node[a] (a4) at (8,-3.6) {Negative binomial: identify $r,p$};
\node[q] (q5) at (0,-4.8) {Event count over fixed exposure?};
\node[a] (a5) at (8,-4.8) {Poisson: find interval mean $\lambda$};
\node[q] (q6) at (0,-6.1) {More information or another model needed};
\foreach \i in {1,2,3,4,5}{\draw[->] (q\i) -- node[above,font=\small]{yes} (a\i);}
\foreach \i/\j in {1/2,2/3,3/4,4/5,5/6}{\draw[->] (q\i) -- node[right,font=\small]{no} (q\j);}
\end{tikzpicture}
\end{center}
If an assumption fails, do not force the scenario into a familiar distribution. State the missing information or explain the unsuitable assumption.
\Needspace{14\baselineskip}
\subsection{Advanced general problem: target versus interval}
Process A consists of independent calls with sale probability 0.4, stopping at the second sale. Process B has isolated arrivals at constant rate 0.5 per hour with independent disjoint-interval counts; observe four hours.

\textbf{Questions/Tasks}
\begin{enumerate}
\item Select models and supports, reject a plausible alternative for each, and calculate the chance A stops on call four and B has at least one arrival.
\end{enumerate}
\subsubsection*{Complete worked solution}
\par\noindent\begin{minipage}{\linewidth}
\textbf{Step 1: Define and select.} $T$ counts calls until sale two, so use negative binomial with $r=2,p=0.4$, support $\{2,3,\ldots\}$ calls. $X$ counts arrivals in four hours, so use Poisson with $\lambda=0.5(4)=2$, support $\{0,1,\ldots\}$ arrivals.
\end{minipage}\par\medskip
\par\noindent\begin{minipage}{\linewidth}
\textbf{Step 2: Check and reject.} A supplies independent constant-probability binary trials, but target two rejects geometric. B supplies Poisson process assumptions; binomial is unsuitable because no fixed number of binary trials is specified.
\end{minipage}\par\medskip
\par\noindent\begin{minipage}{\linewidth}
\textbf{Step 3: Calculate.} $P(T=4)=\binom{4-1}{2-1}p^2(1-p)^2=\binom31(0.4)^2(0.6)^2=0.1728$. Also $P(X\geq1)=1-P(X=0)=1-e^{-\lambda}=1-e^{-2}\approx0.8647$.
\end{minipage}\par\medskip
\par\noindent\begin{minipage}{\linewidth}
\textbf{Step 4: Check and interpret.} A permits SS only if stopping on call two, so fixing the final success on four matters. For B, zero and at least one partition all arrivals. Both answers are probabilities rather than call or arrival counts.
\end{minipage}\par\medskip
\section{Notation reference and practice procedure}
\begin{center}
\begin{tabular}{lp{0.72\linewidth}}
\toprule
Symbol & Meaning and units \\ \midrule
$X$ & Indicator or fixed-exposure success/event count \\
$T$ & Total trials until the specified target success \\
$n$ & Fixed trial count \\
$r$ & Fixed target successes \\
$p$ & Success probability per binary trial \\
$\lambda$ & Expected events in the selected exposure \\
\bottomrule
\end{tabular}
\end{center}
\begin{enumerate}
\item Define the variable and distinguish counts, trial numbers and monetary amounts.
\item Identify what is fixed: one observation, trials, target success or exposure.
\item State support, parameters and units; explicitly define success.
\item Check assumptions and reject at least one plausible alternative.
\item Translate and calculate with a symbolic PMF or complement before substitution.
\item Check the result; if information is insufficient, explain why no unique probability follows.
\end{enumerate}
\Needspace{19\baselineskip}
\section{Foundational contextual practice}\label{sec:t2-c9-foundational}
\Needspace{14\baselineskip}
\subsection{One loan: a binary event}
One loan defaults within a year with probability 0.07; default causes a fixed loss of 500 USD.

\textbf{Questions/Tasks}
\begin{enumerate}
\item Define the indicator, choose and justify the model, reject binomial with ten trials, and find the probability of a 500 USD default loss.
\end{enumerate}
\subsubsection*{Complete worked solution}
\par\noindent\begin{minipage}{\linewidth}
\textbf{Step 1: Define and choose.} $X=1$ for default within a year and 0 otherwise. It is Bernoulli with $p=0.07$, support $\{0,1\}$, no monetary units.
\end{minipage}\par\medskip
\par\noindent\begin{minipage}{\linewidth}
\textbf{Step 2: Justify and reject.} This is one binary observation. No independence assumption is needed. A ten-trial binomial would describe a different variable and ten loans that are not specified.
\end{minipage}\par\medskip
\par\noindent\begin{minipage}{\linewidth}
\textbf{Step 3: Calculate.} $L=500X$ USD, so $P(L=500)=P(X=1)=p^1(1-p)^0=0.07=0.0700$.
\end{minipage}\par\medskip
\par\noindent\begin{minipage}{\linewidth}
\textbf{Step 4: Check and interpret.} The other mass is $1-p=0.93$ and the total is 1. Success denotes an undesirable default; 500 USD is a loss size, not a probability.
\end{minipage}\par\medskip
\secInicio[][sec:t2-c9-foundational]
\Needspace{14\baselineskip}
\subsection{Five independent invoice payments}
Five invoices are independently paid on time, each with probability 0.8. Observe all five.

\textbf{Questions/Tasks}
\begin{enumerate}
\item Define the success count, choose the model and parameters, reject geometric and calculate exactly four on-time payments.
\end{enumerate}
\subsubsection*{Complete worked solution}
\par\noindent\begin{minipage}{\linewidth}
\textbf{Step 1: Define and choose.} $X$ counts on-time payments among five fixed invoices; success means on-time payment. $X\sim\operatorname{Bin}(5,0.8)$ with support $\{0,\ldots,5\}$ payments.
\end{minipage}\par\medskip
\par\noindent\begin{minipage}{\linewidth}
\textbf{Step 2: Justify and reject.} Fixed trials, binary outcomes, independence and constant $p$ are supplied. Geometric would stop at the first on-time payment rather than observe all five.
\end{minipage}\par\medskip
\par\noindent\begin{minipage}{\linewidth}
\textbf{Step 3: Calculate.} $P(X=4)=\binom54p^4(1-p)^{5-4}=5(0.8)^4(0.2)=0.4096$.
\end{minipage}\par\medskip
\par\noindent\begin{minipage}{\linewidth}
\textbf{Step 4: Check and interpret.} The PMF total is $(0.8+0.2)^5=1$. The combination counts which one invoice is late. The answer is a 40.96\% probability, not an expected payment count.
\end{minipage}\par\medskip
\secInicio[][sec:t2-c9-foundational]
\Needspace{19\baselineskip}
\section{Intermediate contextual practice}\label{sec:t2-c9-intermediate}
\Needspace{14\baselineskip}
\subsection{First insurance sale}
Independent offers each sell a policy with probability 0.25. Continue until the first sale.

\textbf{Questions/Tasks}
\begin{enumerate}
\item Define the trial variable, choose and justify the model, reject fixed-trial binomial and calculate first sale after call three.
\end{enumerate}
\subsubsection*{Complete worked solution}
\par\noindent\begin{minipage}{\linewidth}
\textbf{Step 1: Define and choose.} $T$ is total calls up to the first sale; success is a sale. Use geometric with $p=0.25$, support $\{1,2,\ldots\}$ calls.
\end{minipage}\par\medskip
\par\noindent\begin{minipage}{\linewidth}
\textbf{Step 2: Justify and reject.} The independent constant-probability binary process stops at first success. Binomial fixes trials and counts successes, which is a different observation rule.
\end{minipage}\par\medskip
\par\noindent\begin{minipage}{\linewidth}
\textbf{Step 3: Calculate.} $P(T>3)=(1-p)^3=(0.75)^3=0.421875\approx0.4219$.
\end{minipage}\par\medskip
\par\noindent\begin{minipage}{\linewidth}
\textbf{Step 4: Check and interpret.} All first three calls must fail. Its complement $P(T\leq3)=1-0.421875=0.578125$ completes total probability 1. After three excludes sale on call three.
\end{minipage}\par\medskip
\secInicio[][sec:t2-c9-intermediate]
\Needspace{14\baselineskip}
\subsection{Third contract signed on visit five}
Independent visits each sign a contract with probability 0.4. Stop at the third signature.

\textbf{Questions/Tasks}
\begin{enumerate}
\item Define total trials, select parameters and support, reject geometric and find $P(T=5)$. State failures.
\end{enumerate}
\subsubsection*{Complete worked solution}
\par\noindent\begin{minipage}{\linewidth}
\textbf{Step 1: Define and choose.} $T$ counts visits through signature three. Negative binomial has $r=3,p=0.4$ and support $\{3,4,\ldots\}$ visits.
\end{minipage}\par\medskip
\par\noindent\begin{minipage}{\linewidth}
\textbf{Step 2: Justify and reject.} Binary, independent, constant-probability visits stop at target three; geometric targets just one and is unsuitable.
\end{minipage}\par\medskip
\par\noindent\begin{minipage}{\linewidth}
\textbf{Step 3: Calculate.} $P(T=5)=\binom{t-1}{r-1}p^r(1-p)^{t-r}=\binom42(0.4)^3(0.6)^2=0.13824\approx0.1382$.
\end{minipage}\par\medskip
\par\noindent\begin{minipage}{\linewidth}
\textbf{Step 4: Check and interpret.} The first four visits contain two successes; visit five succeeds. There are $F=T-r=5-3=2$ failures. This is not simply three successes anywhere in five visits.
\end{minipage}\par\medskip
\secInicio[][sec:t2-c9-intermediate]
\Needspace{19\baselineskip}
\section{Advanced contextual practice}\label{sec:t2-c9-advanced}
\Needspace{14\baselineskip}
\subsection{Claims in a half-day exposure}
Isolated claims arrive at constant rate 3 per day with independent counts on disjoint intervals. Observe half a day. Each claim costs 200 USD.

\textbf{Questions/Tasks}
\begin{enumerate}
\item Define the count, select and justify the model, reject binomial and calculate the probability of at least 400 USD claim cost.
\end{enumerate}
\subsubsection*{Complete worked solution}
\par\noindent\begin{minipage}{\linewidth}
\textbf{Step 1: Define and choose.} $X$ counts claims in half a day, support $\{0,1,\ldots\}$. Poisson mean is $\lambda=\text{rate}\times\text{exposure}=3(0.5)=1.5$ claims.
\end{minipage}\par\medskip
\par\noindent\begin{minipage}{\linewidth}
\textbf{Step 2: Justify and reject.} Constant rate, independent interval counts and isolated events are supplied. No fixed number of independent binary opportunities is given, so binomial is not justified.
\end{minipage}\par\medskip
\par\noindent\begin{minipage}{\linewidth}
\textbf{Step 3: Calculate.} $L=200X$ USD gives $L\geq400$ when $X\geq2$. Thus $P(X\geq2)=1-[P(X=0)+P(X=1)]=1-e^{-\lambda}(1+\lambda)=1-e^{-1.5}(1+1.5)\approx0.4422$.
\end{minipage}\par\medskip
\par\noindent\begin{minipage}{\linewidth}
\textbf{Step 4: Check and interpret.} The complement includes zero or one claim, cost below 400 USD. The mean 1.5 refers to half a day, not a full day; observed counts remain integers.
\end{minipage}\par\medskip
\secInicio[][sec:t2-c9-advanced]
\Needspace{14\baselineskip}
\subsection{When the information cannot choose a count model}
Five loans each have one-year default probability 0.1. They may share a market shock; their dependence is unspecified. A separate call campaign says only that sales are infrequent, without a numerical sale probability or a stopping rule.

\textbf{Questions/Tasks}
\begin{enumerate}
\item Assess a binomial model for the loan default count, calculate one loan default probability, show that the chance all five default is not uniquely determined, and identify missing campaign information.
\end{enumerate}
\subsubsection*{Complete worked solution}
\par\noindent\begin{minipage}{\linewidth}
\textbf{Step 1: Define and assess.} $X$ counts defaults among five loans, possible values $\{0,\ldots,5\}$. Fixed binary observations are present, but independence is not supplied. Binomial $(5,0.1)$ cannot be asserted. For one specified loan, $I=1$ for default is Bernoulli $(0.1)$ with support $\{0,1\}$.
\end{minipage}\par\medskip
\par\noindent\begin{minipage}{\linewidth}
\textbf{Step 2: Calculate what is known.} $P(I=1)=p^1(1-p)^0=0.1000$. If all five defaults were independent, $P(X=5)=p^5=(0.1)^5=0.00001$. If one common shock caused all five to default together with probability 0.1, then $P(X=5)=0.1$. Both constructions have the supplied individual default probabilities.
\end{minipage}\par\medskip
\par\noindent\begin{minipage}{\linewidth}
\textbf{Step 3: Reject unjustified alternatives.} The different joint probabilities prove no unique aggregate answer follows. Poisson is also unsupported: a fixed five-loan count and possible common shocks do not supply an isolated-event constant-rate process. For the campaign, define whether counting successes in fixed trials or trials to a target, supply $p$ and test independence and constancy before selecting binomial, geometric or negative binomial.
\end{minipage}\par\medskip
\par\noindent\begin{minipage}{\linewidth}
\textbf{Step 4: Check and interpret.} The exact small value $0.00001$ is retained rather than rounded to zero. A single-loan probability can be known while portfolio clustering risk remains unresolved. No numerical campaign probability is justified by the word infrequent.
\end{minipage}\par\medskip
\secInicio[][sec:t2-c9-advanced]
\section{Closing reference: key ideas and common errors}
Choose from the variable and observation rule, then verify assumptions. Bernoulli labels one outcome; binomial counts successes in fixed trials; geometric and negative binomial count total trials to fixed success targets; Poisson counts events in fixed exposure. State support and parameter units, reject a plausible alternative and acknowledge insufficient information. Monetary consequences are separate from event counts.
\secInicio
\end{document}
